\DSource{0142}{53}
\hypertarget{AB01-PDF0142-P01}{}
\DLine{AB01-PDF0142-body-L001}{}{}{\Indent Lineam ME (1) usque ad L producamus, et L cum F coniungamus (2); in triangulo MLF}
\DLine{AB01-PDF0142-body-L002}{}{}{latera proportionalia (3), anguli noti sunt. Arcum AM, ut Ptolemaeus statuit, $120^\circ$ ponamus (4),}
\DLine{AB01-PDF0142-body-L003}{}{}{quae est dupla elongatio Lunae a Sole. Quoniam rationem sinuum ad semidiametrum (5) re-}
\DLine{AB01-PDF0142-body-L004}{}{}{ferimus, angulus EFL (6) $30^\circ$, eiusque complementum FEL $60^\circ$ continebit, iuxta rationem qua}
\DLine{AB01-PDF0142-body-L005}{5}{}{circulus triangulum FEL circumdans $360^\circ$ complectitur. Sinus anguli EFL (7) $30^p$ quoque erit,}
\DLine{AB01-PDF0142-body-L006}{}{}{et sinus anguli FEL circiter $51^p58'$ (8) ex partibus quarum linea EF 60 continet; sed si huic}
\DLine{AB01-PDF0142-body-L007}{}{}{lineae $10^p19'$ tantum (9) tribuantur, linea EL $5^p10'$, linea FL $9^p16'$ (10) circiter fiet. Cum}
\DLine{AB01-PDF0142-body-L008}{}{}{in figura linea EKN epicyclum tetigerit, et locus Lunae in epicyclo fuerit punctum K, erit}
\DLine{AB01-PDF0142-body-L009}{}{}{id maximum quod attingit summa primae inaequalitatis et secundae. Cum autem semidiame-}
\DLine{AB01-PDF0142-body-L010}{10}{}{trus MK epicycli et semidiametrus FM excentrici notae sint, ex ratione FM et FL ratio li-}
\DLine{AB01-PDF0142-body-L011}{}{}{neae ML deprehendetur, eritque tota ML $48^p53'$; de qua, linea EL scilicet $5^p10'$ demptis,}
\DLine{AB01-PDF0142-body-L012}{}{p. 81.}{remanet linea EM a centro [Terrae sive zodiaci] progrediens $43^p43'$. Semidiametrum MK}
\DLine{AB01-PDF0142-body-L013}{}{}{epicycli vidimus $5^p15'$ esse: sed si lineae EM a centro [Terrae vel zodiaci] progredienti $60^p$}
\DLine{AB01-PDF0142-body-L014}{}{}{tribuantur, semidiametrus MK epicycli declivis [ab apogeo excentrici] $7^p12'$ circiter fiet,}
\DLine{AB01-PDF0142-body-L015}{15}{}{et arcus MK super eam insistens $6^\circ54'$ (11) fere continebit. De quo, si tota simplex inaequali-}
\DLine{AB01-PDF0142-body-L016}{}{}{tas, idest $5^\circ1'$, dematur, altera inaequalitas ei adiecta remanebit $1^\circ53'$; et si illi $2^\circ\qfrac{2}{3}$ (12) fiant}
\DLine{AB01-PDF0142-body-L017}{}{}{$60'$, hi $1^\circ53'$ in partes sexagesimas $0^\circ42'38''$ (13) convertentur, quae sub numero 120 in}
\DLine{AB01-PDF0142-body-L018}{}{}{quarta (14) columna reperiuntur. {\itshape Hoc secundum rationem minutorum ad gradum, quae}}
\DLine{AB01-PDF0142-body-L019}{}{}{{\itshape est ratio $0^\circ42'38''$ ad $60'$. Si haec $0^\circ42'38''$ usque ad $60'$ crescunt, ea $1^\circ53'$ fiunt}}
\DLine{AB01-PDF0142-body-L020}{20}{}{{\itshape $2^\circ39'$, quae in quinta} (15) {\itshape columna sub numero 120} (16) {\itshape leguntur} ({\itshape c}).}
\par\noindent\begin{minipage}[t]{0.43\FullSourceWidth}\vspace{0pt}\centering\hypertarget{AB01-PDF0142-F01}{}\includegraphics[width=\linewidth]{AB01-PDF0142-F01.png}\end{minipage}\hfill\begin{minipage}[t]{0.55\FullSourceWidth}\vspace{0pt}\setlength{\GutterOffset}{0.44999999999999996\FullSourceWidth}
\hypertarget{AB01-PDF0142-P02}{}
\DLine{AB01-PDF0142-body-L021}{}{}{\Indent Nunc differentiam inter apogeum verum et apo-}
\DLine{AB01-PDF0142-body-L022}{}{}{geum medium (17), idest arcum HG, cognoscamus.}
\DLine{AB01-PDF0142-body-L023}{}{}{Elongatio Lunae a Sole, iuxta motum medium du-}
\DLine{AB01-PDF0142-body-L024}{}{}{plicatum, sit $90^\circ30'$, ut Ptolemaeus (18) in figura}
\DLine{AB01-PDF0142-body-L025}{25}{}{ex qua id deprehenditur posuit; et motus Lunae in}
\DLine{AB01-PDF0142-body-L026}{}{}{epicyclo a puncto A sit $333^\circ12'$. Ad haec explicanda}
\DLine{AB01-PDF0142-body-L027}{}{}{hunc circulum ponamus.}
\hypertarget{AB01-PDF0142-P03}{}
\DLine{AB01-PDF0142-body-L028}{}{}{\Indent Dicit [auctor]: Circa centrum D diametrumque}
\DLine{AB01-PDF0142-body-L029}{}{}{AC sit excentricus ABC; E ponamus centrum ecli-}
\DLine{AB01-PDF0142-body-L030}{30}{}{pticae, B centrum epicycli MGH. Lineas BM, EBG}
\DLine{AB01-PDF0142-body-L031}{}{}{ducamus, BE usque ad K protrahentes, et K cum}
\DLine{AB01-PDF0142-body-L032}{}{}{D coniungamus, ita ut angulus KDE sit dimidius}
\DLine{AB01-PDF0142-body-L033}{}{}{gradus qui $90^\circ$ excedit; et arcus EK dimidium gra-}\end{minipage}\par
\par\vspace{6pt}\hrule height.25pt\vspace{5pt}
\noindent\begin{minipage}[t]{.48\textwidth}\vspace{0pt}\fontsize{9}{11.5}\selectfont\setlength{\GutterOffset}{0pt}
\hypertarget{AB01-PDF0142-N01}{}
\DLine{AB01-PDF0142-notes_left-L001}{}{}{\Indent (1) Vide figuram pag. 51. --- Iisdem numeris}
\DLine{AB01-PDF0142-notes_left-L002}{}{}{demonstratio legitur in {\itshape Almag.} V, 7 (= ed. Halma,}
\DLine{AB01-PDF0142-notes_left-L003}{}{}{t. I, p. 313--314).}
\hypertarget{AB01-PDF0142-N02}{}
\DLine{AB01-PDF0142-notes_left-L004}{}{}{\Indent (2) FL est perpendicularis ad KL.}
\hypertarget{AB01-PDF0142-N03}{}
\DLine{AB01-PDF0142-notes_left-L005}{}{}{\Indent (3) Interpretandum videtur: « proportiones la-}
\DLine{AB01-PDF0142-notes_left-L006}{}{}{terum sunt notae ».}
\hypertarget{AB01-PDF0142-N04}{}
\DLine{AB01-PDF0142-notes_left-L007}{}{}{\Indent (4) Non arcum AM, sed angulum AEM Ptole-}
\DLine{AB01-PDF0142-notes_left-L008}{}{}{maeus $120^\circ$ complecti statuit; arcus AM subtendit}
\DLine{AB01-PDF0142-notes_left-L009}{}{}{angulum AFM maiorem quam AEM (\Name{Schiapa-}}
\DLine{AB01-PDF0142-notes_left-L010}{}{}{\Name{relli}).}
\hypertarget{AB01-PDF0142-N05}{}
\DLine{AB01-PDF0142-notes_left-L011}{}{}{\Indent (5) Plato addit: « et ad quadrantem circuli ».}
\hypertarget{AB01-PDF0142-N06}{}
\DLine{AB01-PDF0142-notes_left-L012}{}{}{\Indent (6) Cod. et Plato, ut parum infra, AEM.}
\hypertarget{AB01-PDF0142-N07}{}
\DLine{AB01-PDF0142-notes_left-L013}{}{}{\Indent (7) Cod. et Plato, ut supra, AEM.}
\end{minipage}
\hfill\begin{minipage}[t]{.48\textwidth}\vspace{0pt}\fontsize{9}{11.5}\selectfont\setlength{\GutterOffset}{0pt}
\hypertarget{AB01-PDF0142-N08}{}
\DLine{AB01-PDF0142-notes_right-L001}{}{}{\Indent (8) Codex $11^p18'$.}
\hypertarget{AB01-PDF0142-N09}{}
\DLine{AB01-PDF0142-notes_right-L002}{}{}{\Indent (9) Est excentricitas paulo antea reperta.}
\hypertarget{AB01-PDF0142-N10}{}
\DLine{AB01-PDF0142-notes_right-L003}{}{}{\Indent (10) Plato $18^p56'$.}
\hypertarget{AB01-PDF0142-N11}{}
\DLine{AB01-PDF0142-notes_right-L004}{}{}{\Indent (11) Gradus in cod. $5^\circ$.}
\hypertarget{AB01-PDF0142-N12}{}
\DLine{AB01-PDF0142-notes_right-L005}{}{}{\Indent (12) Maximae secundae inaequalitatis, cf. su-}
\DLine{AB01-PDF0142-notes_right-L006}{}{}{pra p. 52.}
\hypertarget{AB01-PDF0142-N13}{}
\DLine{AB01-PDF0142-notes_right-L007}{}{}{\Indent (13) Pro $42'$ Plato semper 45.}
\hypertarget{AB01-PDF0142-N14}{}
\DLine{AB01-PDF0142-notes_right-L008}{}{}{\Indent (14) Vel potius quinta, cfr. supra p. 51, adn. 6}
\DLine{AB01-PDF0142-notes_right-L009}{}{}{et 7; p. 52, adn. 7 et 8.}
\hypertarget{AB01-PDF0142-N15}{}
\DLine{AB01-PDF0142-notes_right-L010}{}{}{\Indent (15) Vel potius sexta.}
\hypertarget{AB01-PDF0142-N16}{}
\DLine{AB01-PDF0142-notes_right-L011}{}{}{\Indent (16) Sed in tabulis sub numero 110 reperiuntur.}
\hypertarget{AB01-PDF0142-N17}{}
\DLine{AB01-PDF0142-notes_right-L012}{}{}{\Indent (17) In tabulis {\itshape aequatio anomaliae} appellatur.}
\hypertarget{AB01-PDF0142-N18}{}
\DLine{AB01-PDF0142-notes_right-L013}{}{}{\Indent (18) {\itshape Almag.} V, 6 (ed. Halma t. I, p. 308--311).}
\end{minipage}
